\documentclass[11pt]{article}

\usepackage[margin=1in]{geometry}
\usepackage{amsmath}
\usepackage{amssymb}
\usepackage{mathtools}
\usepackage{titlesec}
\usepackage{parskip}
\usepackage[hidelinks]{hyperref}

\titleformat{\section}{\normalfont\Large\bfseries}{\thesection}{1em}{}
\renewcommand{\arraystretch}{1.3}

\title{Linear Stability To Check}
\date{2026-08-17}
\author{}

\begin{document}
\maketitle

\section*{Step 0: The model}

\[
\partial_t n = n\left(c_{\text{ext}}(1-l)\,r_{\text{ext}} + c_\times r_\times - m\right) + D_n \nabla^2 n
\]

\[
\partial_t r_{\text{ext}} = K - \kappa\, r_{\text{ext}} - c_{\text{ext}}\, r_{\text{ext}}\, n + D_{\text{ext}} \nabla^2 r_{\text{ext}}
\]

\[
\partial_t r_\times = c_{\text{ext}}\, l\, r_{\text{ext}}\, n - \kappa\, r_\times - c_\times r_\times n + D_\times \nabla^2 r_\times
\]

The story: cells consume $r_{\text{ext}}$ at rate $c_{\text{ext}} r_{\text{ext}} n$. A fraction $l$ of that flux is leaked as the byproduct $r_\times$; the remaining fraction $(1-l)$ becomes growth. The leaked byproduct is then re-consumed by the same population at rate $c_\times r_\times n$, also contributing to growth.

\bigskip\hrule\bigskip

\section*{Step 1: Homogeneous steady state}

Set all time derivatives and all Laplacians to zero.

\textbf{From the $r_{\text{ext}}$ equation:}

\[
K = \kappa r_{\text{ext}}^* + c_{\text{ext}} r_{\text{ext}}^* n^* \quad\Longrightarrow\quad r_{\text{ext}}^* = \frac{K}{\kappa + c_{\text{ext}} n^*}
\]

\textbf{From the $r_\times$ equation:}

\[
c_{\text{ext}} l\, r_{\text{ext}}^* n^* = \kappa r_\times^* + c_\times r_\times^* n^* \quad\Longrightarrow\quad r_\times^* = \frac{c_{\text{ext}} l\, r_{\text{ext}}^* n^*}{\kappa + c_\times n^*}
\]

These two denominators will appear everywhere, so name them now:

\[
\boxed{\;u \equiv \kappa + c_{\text{ext}} n^*, \qquad v \equiv \kappa + c_\times n^*\;}
\]

Each is a total removal rate: $u$ is the rate at which $r_{\text{ext}}$ leaves the system (dilution $\kappa$ plus consumption $c_{\text{ext}}n^*$), and $v$ likewise for $r_\times$. In these terms

\[
r_{\text{ext}}^* = \frac{K}{u}, \qquad r_\times^* = \frac{c_{\text{ext}} l K n^*}{uv}
\]

\textbf{From the $n$ equation}, the bracket must vanish (taking $n^*\neq 0$):

\[
c_{\text{ext}}(1-l)\, r_{\text{ext}}^* + c_\times r_\times^* = m
\]

This is the equation that fixes $n^*$. We won't need to solve it explicitly --- but we \emph{will} use the fact that it holds, repeatedly.

\bigskip\hrule\bigskip

\section*{Step 2: Perturb}

Write

\[
n = n^* + \nu\, e^{i\mathbf{k}\cdot\mathbf{x} + \lambda t}, \qquad r_{\text{ext}} = r_{\text{ext}}^* + \rho\, e^{i\mathbf{k}\cdot\mathbf{x}+\lambda t}, \qquad r_\times = r_\times^* + \sigma\, e^{i\mathbf{k}\cdot\mathbf{x}+\lambda t}
\]

with $\nu, \rho, \sigma$ infinitesimal. Since $\nabla^2 e^{i\mathbf{k}\cdot\mathbf{x}} = -k^2 e^{i\mathbf{k}\cdot\mathbf{x}}$, every Laplacian becomes multiplication by $-k^2$.

\textbf{Quasi-steady-state assumption:} the resources equilibrate fast, so we set $\partial_t \rho = \partial_t\sigma = 0$ while keeping their diffusion. This is the key simplification: it lets us solve for $\rho$ and $\sigma$ in terms of $\nu$, and reduces a $3\times3$ eigenvalue problem to a single scalar equation.

\bigskip\hrule\bigskip

\section*{Step 3: Solve for $\rho$ (the external resource response)}

Linearise the $r_{\text{ext}}$ equation. The only nonlinear term is $-c_{\text{ext}} r_{\text{ext}} n$, which varies as $-c_{\text{ext}}(r_{\text{ext}}^*\nu + n^*\rho)$:

\[
0 = -\kappa\rho - c_{\text{ext}}\left(r_{\text{ext}}^*\nu + n^*\rho\right) - D_{\text{ext}}k^2\rho
\]

Collect the $\rho$ terms:

\[
0 = -\underbrace{\left(\kappa + c_{\text{ext}}n^*\right)}_{=\,u}\rho - D_{\text{ext}}k^2\rho - c_{\text{ext}}r_{\text{ext}}^*\nu
\]

\[
\boxed{\;\rho = -\frac{c_{\text{ext}}\, r_{\text{ext}}^*}{u + D_{\text{ext}}k^2}\;\nu\;}
\]

\textbf{Read it:} $\rho$ has the \emph{opposite} sign to $\nu$. More cells means less external resource. This is the inhibitory channel. Its magnitude is suppressed at large $k$ --- because rapid diffusion smooths out any local depletion, a short-wavelength bloom cannot dig itself a resource hole.

\bigskip\hrule\bigskip

\section*{Step 4: Solve for $\sigma$ (the byproduct response)}

Linearise the $r_\times$ equation. Two nonlinear terms: $+c_{\text{ext}}l\, r_{\text{ext}}n \to c_{\text{ext}}l(r_{\text{ext}}^*\nu + n^*\rho)$, and $-c_\times r_\times n \to -c_\times(r_\times^*\nu + n^*\sigma)$:

\[
0 = c_{\text{ext}}l\left(r_{\text{ext}}^*\nu + n^*\rho\right) - \kappa\sigma - c_\times\left(r_\times^*\nu + n^*\sigma\right) - D_\times k^2 \sigma
\]

Collect the $\sigma$ terms, which give $-(\kappa + c_\times n^*)\sigma - D_\times k^2\sigma = -(v + D_\times k^2)\sigma$:

\[
\sigma = \frac{c_{\text{ext}}l\, r_{\text{ext}}^*\nu \;-\; c_\times r_\times^* \nu \;+\; c_{\text{ext}}l\, n^*\rho}{v + D_\times k^2}
\]

Now handle the first two terms. Using $r_\times^* = c_{\text{ext}}lKn^*/(uv)$ and $r_{\text{ext}}^*=K/u$:

\[
c_{\text{ext}}l\,r_{\text{ext}}^* - c_\times r_\times^* = \frac{c_{\text{ext}}lK}{u} - \frac{c_\times c_{\text{ext}}lKn^*}{uv} = \frac{c_{\text{ext}}lK}{u}\left(1 - \frac{c_\times n^*}{v}\right) = \frac{c_{\text{ext}}lK}{u}\cdot\frac{\kappa}{v}
\]

using $v - c_\times n^* = \kappa$. So

\[
\boxed{\;\sigma = \frac{1}{v+D_\times k^2}\left[\frac{c_{\text{ext}}lK\kappa}{uv}\,\nu \;+\; c_{\text{ext}}l\,n^*\rho\right]\;}
\]

\textbf{Read it:} two contributions.
\begin{itemize}
  \item The first is \textbf{positive}: more cells leak more byproduct. This is the activating channel.
  \item The second inherits $\rho$, which is negative. More cells deplete $r_{\text{ext}}$, so \emph{less} gets leaked into $r_\times$. This is an inhibitory contribution travelling through the byproduct.
\end{itemize}

That second piece is the one that will matter most, so let's track it carefully.

\bigskip\hrule\bigskip

\section*{Step 5: The $n$ equation}

Linearise. The perturbation $\nu$ multiplies the bracket --- but \textbf{the bracket vanishes at the fixed point} (Step 1). So that term dies, and only $n^*$ times the bracket's perturbation survives:

\[
\lambda\nu = n^*\left[c_{\text{ext}}(1-l)\rho + c_\times \sigma\right] - D_n k^2 \nu
\]

This is the crucial structural fact: \textbf{$n$ has no autonomous linear dynamics.} Its entire growth rate is inherited from the resource perturbations. Everything downstream follows from this.

\bigskip\hrule\bigskip

\section*{Step 6: Assemble}

Substitute $\sigma$ from Step 4:

\[
\lambda = -D_nk^2 + \frac{n^*}{\nu}\left[c_{\text{ext}}(1-l)\rho + \frac{c_\times}{v+D_\times k^2}\left(\frac{c_{\text{ext}}lK\kappa}{uv}\nu + c_{\text{ext}}l\,n^*\rho\right)\right]
\]

Now substitute $\rho = -\dfrac{c_{\text{ext}}K/u}{u+D_{\text{ext}}k^2}\nu$ in the two places it appears, and divide out $\nu$:

\[
\lambda = -D_nk^2 \;\underbrace{-\; \frac{n^*c_{\text{ext}}^2(1-l)K}{u\left(u+D_{\text{ext}}k^2\right)}}_{\text{direct competition}} \;\underbrace{+\; \frac{n^*c_\times c_{\text{ext}}lK\kappa}{uv\left(v+D_\times k^2\right)}}_{\text{cross-feeding gain}} \;\underbrace{-\; \frac{n^{*2}c_\times c_{\text{ext}}^2 lK}{u\left(u+D_{\text{ext}}k^2\right)\left(v+D_\times k^2\right)}}_{\text{indirect: depletion} \to \text{less leakage}}
\]

\bigskip\hrule\bigskip

\section*{Step 7: Compact form}

Factor each denominator to expose its structure. Define the two \textbf{screening lengths}

\[
\ell_{\text{ext}} = \sqrt{\frac{D_{\text{ext}}}{u}}, \qquad \ell_\times = \sqrt{\frac{D_\times}{v}}
\]

so that $u + D_{\text{ext}}k^2 = u\left(1+\ell_{\text{ext}}^2k^2\right)$ and $v+D_\times k^2 = v\left(1+\ell_\times^2k^2\right)$. Then, defining

\[
A = \frac{n^*c_\times c_{\text{ext}}lK\kappa}{uv^2}, \qquad I_1 = \frac{n^*c_{\text{ext}}^2(1-l)K}{u^2}, \qquad I_2 = \frac{n^{*2}c_\times c_{\text{ext}}^2lK}{u^2v}
\]

we arrive at

\[
\boxed{\;\lambda(k) = -D_nk^2 \;+\; \frac{A}{1+\ell_\times^2k^2} \;-\; \frac{I_1}{1+\ell_{\text{ext}}^2k^2} \;-\; \frac{I_2}{\left(1+\ell_{\text{ext}}^2k^2\right)\left(1+\ell_\times^2k^2\right)}\;}
\]

All three of $A, I_1, I_2$ are \textbf{positive}. This is an activator--inhibitor dispersion relation: one gain term, two loss terms.

\bigskip\hrule\bigskip

\section*{Step 8: Sanity checks}

\textbf{At $k=0$:} all Lorentzians equal $1$, so $\lambda(0) = A - I_1 - I_2$. This must be negative for the homogeneous state to be stable --- and indeed it equals $n^*G'(n^*)$, the well-mixed eigenvalue. Net feedback is inhibitory in bulk.

\textbf{At $k\to\infty$ with $D_n=0$:} every term vanishes, $\lambda\to 0$. This is Step 5's consequence: screen out the resources and $n$ has nothing left to do.

\bigskip\hrule\bigskip

\section*{Step 9: What drives the instability}

Since $\lambda(0)<0$ and $\lambda(\infty)=0$ (for $D_n=0$), instability is decided by \emph{how} zero is approached. Expand at large $k$:

\[
\frac{A}{1+\ell_\times^2k^2} \sim \frac{A}{\ell_\times^2k^2}, \qquad \frac{I_1}{1+\ell_{\text{ext}}^2k^2}\sim\frac{I_1}{\ell_{\text{ext}}^2k^2}, \qquad \frac{I_2}{(1+\ell_{\text{ext}}^2k^2)(1+\ell_\times^2k^2)} \sim \frac{I_2}{\ell_{\text{ext}}^2\ell_\times^2 k^4}
\]

\textbf{The key observation:} $I_2$ decays as $k^{-4}$, one power faster than the other two. It is a \emph{double} Lorentzian, because that inhibitory signal must pass through \textbf{two} screened resource fields in sequence --- depletion of $r_{\text{ext}}$, propagated into $r_\times$. In real space this means its range is $\sqrt{\ell_{\text{ext}}^2 + \ell_\times^2}$, strictly longer than the activator's $\ell_\times$, \textbf{for any diffusivities whatsoever}.

So the separation of scales here is \textbf{topological, not diffusive} --- it comes from path length through the metabolic chain, not from $D_{\text{ext}}$ exceeding $D_\times$. Dropping $I_2$ from the large-$k$ balance:

\[
\lambda(k) \approx \frac{1}{k^2}\left(\frac{A}{\ell_\times^2} - \frac{I_1}{\ell_{\text{ext}}^2}\right) = \frac{1}{k^2}\left(\frac{Av}{D_\times} - \frac{I_1 u}{D_{\text{ext}}}\right)
\]

Instability requires this to be positive. Substituting the definitions, $Av/D_\times = \dfrac{n^*c_\times c_{\text{ext}}lK\kappa}{uvD_\times}$ and $I_1u/D_{\text{ext}} = \dfrac{n^*c_{\text{ext}}^2(1-l)K}{uD_{\text{ext}}}$, so the condition is

\[
\boxed{\;\frac{1}{D_\times}\cdot\frac{c_\times l\,\kappa}{v} \;>\; \frac{1}{D_{\text{ext}}}\cdot c_{\text{ext}}(1-l)\;}
\]

For $D_{\text{ext}} = D_\times$ this reduces to the reaction-only condition $\dfrac{c_\times l\kappa}{\kappa+c_\times n^*} > c_{\text{ext}}(1-l)$ --- strong leakage and weak direct competition --- with no diffusivity requirement, which is the case you checked against $Kc_{\text{ext}}/(m\kappa) < 1/(1-l)$.

\textbf{Two necessary ingredients, distinct from each other:}

\begin{enumerate}
  \item \textbf{Sign:} the above inequality, which asks whether the net feedback flips positive once the resources are screened. Mostly a statement about reaction kinetics.
  \item \textbf{Scale:} at least one $D>0$, so that screening actually happens as $k$ grows. If $D_{\text{ext}}=D_\times=0$ there is no $k$-dependence at all and $\lambda(k)\equiv\lambda_0<0$ everywhere.
\end{enumerate}

And $D_n$ must be small: it is the only term that \emph{damps} rather than merely screening, and a large enough $D_n$ kills the band regardless of everything else. With $D_n>0$ the unstable band is finite and a genuine wavelength is selected; with $D_n = 0$ the band extends to $k\to\infty$ and the selected scale must come from elsewhere.

\end{document}
